\section{Quick overview about Knowledge Graph}

``Knowledge graphs triple LLM accuracy on enterprise data. But here's what nobody tells you upfront: every knowledge graph converges on the same patterns, the same solutions.''
(\href{https://bryon.io/why-rdf-is-the-natural-knowledge-layer-for-ai-systems-a5fd0b43d4c5}{the series “LLMs Need Knowledge Graphs. Use RDF or End Up Rebuilding It.”}).

It means that when organizations build knowledge graphs, they often end up discovering that they need the same fundamental structures, regardless of their original technology choice.

For example, a company may start with:
\begin{itemize}
  \item a graph database like Neo4j
  \item a property graph model
  \item a custom JSON graph
  \item a relational database with relationships
  \item an ontology system
\end{itemize}

\textbf{But over time, they keep needing the same capabilities} (Table~\ref{tab:kg-patterns}).
\begin{tabularx}{\textwidth}{@{}>{\raggedright\arraybackslash}p{3.8cm}X@{}}
  \caption{Patterns every knowledge graph converges on}
  \label{tab:kg-patterns}                            \\

  \toprule
  \textbf{Pattern} & \textbf{Description \& Example} \\
  \midrule
  \endfirsthead

  \toprule
  \textbf{Pattern} & \textbf{Description \& Example} \\
  \midrule
  \endhead

  \bottomrule
  \endfoot


  \textbf{Pattern 1:}                                \\
  Everything becomes an entity
                   &
  Companies realize they need to represent identifiable objects such as:

  \medskip

  \begin{tabular}[t]{@{}l@{}}
    \texttt{Customer} \\
    \texttt{Employee} \\
    \texttt{Product}  \\
    \texttt{Contract} \\
    \texttt{Supplier} \\
    \texttt{Location} \\
    \texttt{Document}
  \end{tabular}
  \\

  \addlinespace


  \textbf{Pattern 2:}                                \\
  Everything becomes a relationship
                   &
  They need to express connections between entities:

  \medskip

  \begin{tabular}[t]{@{}l@{}}
    \texttt{Customer} $\rightarrow$ \texttt{purchased} $\rightarrow$ \texttt{Product}     \\
    \texttt{Employee} $\rightarrow$ \texttt{works\_for} $\rightarrow$ \texttt{Department} \\
    \texttt{Contract} $\rightarrow$ \texttt{applies\_to} $\rightarrow$ \texttt{Product}   \\
    \texttt{Disease} $\rightarrow$ \texttt{treated\_by} $\rightarrow$ \texttt{Drug}
  \end{tabular}
  \\

  \addlinespace


  \textbf{Pattern 3:}                                \\
  Shared meaning (semantics)
                   &
  Different systems often use different names for the same concept:

  \medskip

  \begin{tabular}[t]{@{}l@{}}
    CRM: \texttt{"Client"}   \\
    ERP: \texttt{"Customer"} \\
    Support system: \texttt{"Account"}
  \end{tabular}

  \medskip

  A knowledge graph reconciles these concepts as:

  \medskip

  \begin{center}
    \texttt{Client = Customer = Account}
  \end{center}

  This is where \emph{ontologies} come into play.
  \\

  \addlinespace


  \textbf{Pattern 4:}                                \\
  Reasoning
                   &
  Organizations eventually ask questions such as:

  \medskip

  \emph{``Which suppliers are affected by this regulation?''}

  \medskip

  Although the supplier may not be directly linked to the regulation, the graph can infer the relationship through intermediate entities:

  \medskip

  \begin{tabular}[t]{@{}l@{}}
    \texttt{Supplier}                  \\
    $\downarrow$ \texttt{supplies}     \\
    \texttt{Product}                   \\
    $\downarrow$ \texttt{contains}     \\
    \texttt{Component}                 \\
    $\downarrow$ \texttt{affected\_by} \\
    \texttt{Regulation}
  \end{tabular}
  \\


  \bottomrule
\end{tabularx}